% =========================================================================
\section{El campo eléctrico}\label{sec:campo}
% =========================================================================
La ley de Coulomb dice qué fuerza sufre una carga, pero no cómo repartir el trabajo entre «quien crea» y
«quien siente». Para separar ambos papeles, llamamos \emph{carga fuente} $\qp$ a la que crea la perturbación,
situada en el \emph{punto fuente} $\vc{r}\,'$, y \emph{carga de prueba} $q$ (que tomamos positiva) a la que
experimenta la fuerza, situada en el \emph{punto campo} $\vc{r}$. La carga de prueba no se incluye entre las
fuentes: ninguna carga ejerce fuerza sobre sí misma (\cref{fig:campo}a).

% -------------------------------------------------------------------------
\subsection{Definición del campo eléctrico}\label{sec:def-campo}
% -------------------------------------------------------------------------
Dividiendo la fuerza de Coulomb por la carga de prueba, desaparece toda dependencia de $q$ y queda algo asociado
únicamente a las fuentes.

\begin{definicion}[Campo eléctrico]\label{def:campo}
El \emph{campo eléctrico} en el punto $\vc{r}$ es la fuerza por unidad de carga que experimentaría una carga de prueba
en ese punto:
\begin{equation}\label{eq:campo-def}
  \vE(\vc{r})=\lim_{q\to0}\frac{\vF_q}{q}.
\end{equation}
Para una carga fuente puntual $\qp$ situada en $\vc{r}\,'$,
\begin{equation}\label{eq:campo-puntual}
  \boxed{\;\vE(\vc{r})=\frac{1}{4\pi\eo}\,\frac{\qp}{R^{2}}\,\uR
  =\frac{1}{4\pi\eo}\,\frac{\qp}{\norm{\vc{R}}^{3}}\,\vc{R},\qquad\vc{R}=\vc{r}-\vc{r}\,'\;}
\end{equation}
Su unidad en el SI es el N/C, aunque se usa con más frecuencia el V/m (\cref{sec:potencial-signif}).
\end{definicion}

\begin{observacion}[Por qué el límite $q\to0$]\label{obs:limite}
Con $\vE=\vF_q/q$ bastaría para una carga fuente «rígida», pero si la fuente es, por ejemplo, una esfera
conductora, acercar una carga de prueba apreciable redistribuye la carga $\qp$ sobre la esfera y el campo que
se quiere medir cambia, haciéndose dependiente de la propia carga de prueba. El límite $q\to0$ corrige esto: la
carga de prueba debe ser tan pequeña que no perturbe la distribución que crea el campo. Es una idealización, pues
existe una carga mínima, la del electrón.
\end{observacion}

Con esta definición, el campo es una función $\vE:\R^{3}\to\R^{3}$ (un \emph{campo vectorial}) y la fuerza
sobre cualquier carga $q$ colocada en $\vc{r}$ se obtiene como $\vF_q=q\,\vE(\vc{r})$.

% Figura: carga fuente, punto campo y distribuciones continuas de carga.
\begin{figure}[ht]
\centering
\begin{minipage}[t]{0.40\textwidth}\centering
\tdplotsetmaincoords{72}{125}
\begin{tikzpicture}[tdplot_main_coords, scale=0.95, baseline=(current bounding box.north)]
  \draw[eje] (0,0,0) -- (3.0,0,0) node[below left] {\small $x$};
  \draw[eje] (0,0,0) -- (0,3.2,0) node[right] {\small $y$};
  \draw[eje] (0,0,0) -- (0,0,3.0) node[above] {\small $z$};
  \coordinate (S) at (2.4,0.3,0.4);
  \coordinate (P) at (0.4,2.6,2.2);
  \draw[->, azulclaro!80!azulacero, line width=0.7pt] (0,0,0) -- (S) node[pos=0.5, below, font=\small, text=azulprof] {$\vc{r}\,'$};
  \draw[->, azulclaro!80!azulacero, line width=0.7pt] (0,0,0) -- (P) node[pos=0.45, left, font=\small, text=azulprof] {$\vc{r}$};
  \draw[vect, azulprof, line width=1.2pt] (S) -- (P) node[pos=0.55, right=2pt, font=\small, text=azulprof] {$\vc{R}$};
  \draw[vectrojo, line width=1.2pt] (P) -- ($(P)+(-0.45,0.7,0.65)$) node[above, font=\small, text=acento] {$\vE(\vc{r})$};
  \node[cargapos] at (S) {$+$};
  \node[cargapru] at (P) {$q$};
  \node[font=\small, anchor=south west] at ($(S)+(0.1,0.1,0.1)$) {$\qp$};
\end{tikzpicture}\\[2pt]{\small\scshape (a)}
\end{minipage}\hfill
\begin{minipage}[t]{0.56\textwidth}\centering
\begin{tikzpicture}[baseline=(current bounding box.north), x=1cm, y=1cm]
  % (b) volumen
  \begin{scope}[shift={(0,0)}]
    \draw[azulprof, fill=azulfondo, line width=0.6pt] (0,0.2) .. controls (-0.2,1.2) and (0.4,1.9) .. (1.1,1.8) .. controls (1.9,1.9) and (2.2,0.9) .. (1.7,0.3) .. controls (1.2,-0.3) and (0.4,-0.2) .. cycle;
    \fill[azulacero] (1.0,1.0) circle (1.3pt);
    \draw[azulacero, fill=azulclaro!60, line width=0.4pt] (0.95,0.98) rectangle (1.12,1.13);
    \node[font=\scriptsize, anchor=south east] at (0.95,1.1) {$\dif V\,'$};
    \node[font=\scriptsize, text=azulprof] at (0.4,0.35) {$V'$};
    \node[font=\scriptsize, anchor=north] at (1.0,-0.5) {$\dif q'=\rho\,\dVp$};
  \end{scope}
  % (c) superficie
  \begin{scope}[shift={(2.8,0)}]
    \draw[azulprof, fill=azulfondo, line width=0.6pt] (0,0.1) -- (0.3,1.5) .. controls (1.0,1.9) and (1.5,1.4) .. (2.0,1.7) -- (1.8,0.4) .. controls (1.2,0.0) and (0.8,0.4) .. cycle;
    \draw[azulacero, fill=azulclaro!60, line width=0.4pt] (0.9,0.8) -- (1.15,0.85) -- (1.1,1.1) -- (0.85,1.05) -- cycle;
    \node[font=\scriptsize, anchor=south east] at (0.85,1.1) {$\dSp$};
    \node[font=\scriptsize, text=azulprof] at (0.4,0.45) {$S'$};
    \node[font=\scriptsize, anchor=north] at (1.0,-0.5) {$\dif q'=\sigma\,\dSp$};
  \end{scope}
  % (d) línea
  \begin{scope}[shift={(5.6,0)}]
    \draw[azulprof, line width=1.1pt] (0,0.3) .. controls (0.5,1.5) and (1.0,-0.4) .. (1.6,1.0) .. controls (1.9,1.5) and (2.0,1.6) .. (2.1,1.7);
    \draw[azulacero, line width=2.8pt] (0.78,0.64) -- (0.99,0.62);
    \node[font=\scriptsize, anchor=south] at (0.85,0.95) {$\dlp$};
    \node[font=\scriptsize, text=azulprof] at (0.05,0.7) {$l'$};
    \node[font=\scriptsize, anchor=north] at (1.0,-0.5) {$\dif q'=\lambda\,\dlp$};
  \end{scope}
\end{tikzpicture}\\[2pt]{\small\scshape (b)}
\end{minipage}
\caption{(a)~Carga fuente $\qp$ en $\vc{r}\,'$ y carga de prueba $q$ en el punto campo $\vc{r}$;
$\vc{R}=\vc{r}-\vc{r}\,'$ va de la fuente al punto campo, y el campo $\vE(\vc{r})$ de una
fuente positiva apunta en su dirección. (b)~Distribuciones continuas: de volumen $V'$
($\dif q'=\rho\,\dVp$), de superficie $S'$ ($\dif q'=\sigma\,\dSp$) y de línea $l'$ ($\dif q'=\lambda\,\dlp$);
cada elemento actúa como una carga puntual $\dif q'$ en su propia posición $\vc{r}\,'$.}
\label{fig:campo}
\end{figure}


% -------------------------------------------------------------------------
\subsection{Campo de una distribución discreta de cargas}\label{sec:campo-discreto}
% -------------------------------------------------------------------------
\begin{proposicion}[Campo de un sistema de cargas puntuales]\label{prop:campo-discreto}
El campo eléctrico creado en $\vc{r}$ por las cargas puntuales $\qp_1,\dots,\qp_N$ situadas en
$\vc{r}\,'_1,\dots,\vc{r}\,'_N$ es
\begin{equation}\label{eq:campo-discreto}
  \vE(\vc{r})=\frac{1}{4\pi\eo}\sum_{i=1}^{N}\frac{\qp_i}{R_i^{2}}\,\uh{R_i}
  =\frac{1}{4\pi\eo}\sum_{i=1}^{N}\frac{\qp_i}{\norm{\vc{R}_i}^{3}}\,\vc{R}_i,\qquad
  \vc{R}_i=\vc{r}-\vc{r}\,'_i .
\end{equation}
\end{proposicion}
\begin{proof}
Por el principio de superposición (\cref{eq:superposicion}), la fuerza sobre una carga de prueba $q$ es la
suma de las fuerzas de cada $\qp_i$. Dividiendo por $q$ y tomando el límite $q\to0$ se obtiene
\cref{eq:campo-discreto}.
\end{proof}

% -------------------------------------------------------------------------
\subsection{Distribuciones continuas de carga}\label{sec:campo-continuo}
% -------------------------------------------------------------------------
Cuando las cargas están muy próximas entre sí en comparación con las demás distancias de interés, resulta conveniente
tratarlas como un continuo y definir densidades de carga.

\begin{definicion}[Densidades de carga]\label{def:densidades}
Sea $\dif q'$ la carga contenida en un elemento de volumen $\dVp$ (respectivamente, de superficie $\dSp$ o de
línea $\dlp$) situado en $\vc{r}\,'$. Se definen las densidades \emph{volúmica}, \emph{superficial} y \emph{lineal}:
\begin{align}
  \rho(\vc{r}\,')&=\frac{\dif q'}{\dVp}, & Q&=\int_{V'}\rho(\vc{r}\,')\,\dVp,\label{eq:rho}\\
  \sigma(\vc{r}\,')&=\frac{\dif q'}{\dSp}, & Q&=\int_{S'}\sigma(\vc{r}\,')\,\dSp,\label{eq:sigma}\\
  \lambda(\vc{r}\,')&=\frac{\dif q'}{\dlp}, & Q&=\int_{l'}\lambda(\vc{r}\,')\,\dlp.\label{eq:lambda}
\end{align}
\end{definicion}

\begin{observacion}[Elemento macroscópico]\label{obs:macro}
En \cref{eq:rho}, $\dVp$ debe ser un volumen \emph{pequeño a escala macroscópica}, de modo que permita seguir las
variaciones de la distribución y tratar $\dif q'$ como una carga puntual, pero \emph{grande a escala microscópica},
de modo que contenga muchos átomos o moléculas. Si $\dVp$ fuese menor que un átomo, en casi todas las posiciones no
habría carga y $\rho$ sería nula, salvo cuando $\dVp$ incluyese un electrón, y esto produciría fluctuaciones enormes en
$\rho$ que lo harían inservible como concepto. Así, $\rho$ representa un valor \emph{medio} que varía de forma suave
y que podemos considerar una función continua de la posición.
\end{observacion}

\begin{proposicion}[Campo de una distribución continua]\label{prop:campo-continuo}
El campo eléctrico creado en $\vc{r}$ por una distribución continua de carga es
\begin{align}
  \vE(\vc{r})&=\frac{1}{4\pi\eo}\int_{V'}\frac{\dif q'}{R^{2}}\,\uR
  =\frac{1}{4\pi\eo}\int_{V'}\frac{\rho(\vc{r}\,')}{\norm{\vc{R}}^{3}}\,\vc{R}\;\dVp,\label{eq:campo-vol}\\
  \vE(\vc{r})&=\frac{1}{4\pi\eo}\int_{S'}\frac{\sigma(\vc{r}\,')}{\norm{\vc{R}}^{3}}\,\vc{R}\;\dSp,\qquad
  \vE(\vc{r})=\frac{1}{4\pi\eo}\int_{l'}\frac{\lambda(\vc{r}\,')}{\norm{\vc{R}}^{3}}\,\vc{R}\;\dlp,\label{eq:campo-sup-lin}
\end{align}
con $\vc{R}=\vc{r}-\vc{r}\,'$.
\end{proposicion}
\begin{proof}
Cada elemento $\dif q'$ de la distribución se comporta como una carga puntual en $\vc{r}\,'$, que crea un campo
elemental $\dif\vE=\frac{1}{4\pi\eo}\frac{\dif q'}{R^{3}}\vc{R}$ por \cref{eq:campo-puntual}. El principio de
superposición, con la suma sobre elementos convertida en integral (el límite de las sumas de
Riemann de \cref{eq:campo-discreto}), da \cref{eq:campo-vol,eq:campo-sup-lin}, sin más que expresar
$\dif q'$ como $\rho\,\dVp$, $\sigma\,\dSp$ o $\lambda\,\dlp$.
\end{proof}

\begin{observacion}[Cuidado con las variables]\label{obs:variables}
Dentro de la integral, $\vc{r}\,'$ y, por tanto, $\vc{R}$ \emph{varían} para cada elemento de la distribución;
$\vc{r}$, en cambio, es el punto campo y permanece fijo. Si la distribución tiene la simetría adecuada, $\vE$ puede
calcularse por integración directa de las expresiones anteriores (veremos varios casos en los
seminarios, \cref{sec:seminarios}), pero esto no es lo habitual: en la mayoría de los casos la integral es difícil, y
obtendremos mucho más partido de las ecuaciones diferenciales que se deducen en las secciones siguientes.
\end{observacion}
